\cotefile{157-4}
\scanpage{180}
\mspage{54}
% archive sheet number: 179

\unsure{prouver} que $M$ est stable par $\varinjlim$ grandes devant $\unsure{\pi|\Pi}$ sans plus, et je doute maintenant que ça résulte des seules hypothèses faites ici, qui ne concernent que les systèmes inductifs où les morphismes de transition sont mono.

Voici, je crois, un contre-exemple. Prenons

\[
\begin{aligned}
 M={}&\text{ensemble des objets de $\mathcal M$ isomorphes}\\
 &\text{à un objet \unsure{fixe|fixé} $\deleted{\unsure{X}}X_0\in\mathrm{Filt}^{\Pi}_{\deleted{\unsure{\pi}}<}$,}\\
 &\text{tel que $X_0$ n'ait pas d'\unsure{endomorphismes|endomorphisme}}\\
 &\text{qui soit mono que l'identité.}
\end{aligned}\tag*{(5.4.)}
\]

Alors la condition 5.3. est satisfaite. Car \deleted{\unsure{a)}}
a) si $X\in M$, i.e. $X\simeq X_0$, donc $X\in\mathrm{Filt}^{\Pi}_{\deleted{\unsure{\pi|<}}}$, \deleted{\unsure{donc}}
$\mathrm{Ss}_{\pi}(X)$ admet \deleted{\unsure{un}} \inserted{$X$ comme} objet final et $J$, \deleted{étant} cofinal, contient \deleted{$X$,} \inserted{cet \unsure{élément|objet}}, donc $J_M$ (qui le contient aussi) est cofinal. Et b) si $J_M$ est cofinal, \deleted{\illegible{1}} l'hypothèse (5.4.) implique que les morphismes \deleted{de} transition \deleted{\unsure{dans}} entre les $Z_\alpha$ ($\alpha\in J_M$) est formé d'isomorphismes, \deleted{donc} la limite inductive $X$ des $Z_\alpha$ en question est isom. à chacun d'eux, \deleted{donc} $X\deleted{\in M.}\simeq X_0$ i.e. $X\in M$.

\unsure{Pourtant|Par contre}, $M$ n'est $\pi$-accessible que s'il est stable par $\varinjlim$ grandes devant $\pi$, ce qui implique qu'il est stable par \uline{facteurs directs}. \deleted{\unsure{Donc Mais l'hypothèse}}

